\subsection{Continuum and Constitutive Assumptions}
\label{subsec:continuum-constitutive-assumptions}

This section moves from the macroscopic scale assumption, 
to continuum field variables, to balance-law vocabulary,
and finally to Newtonian versus generalized-Newtonian stress closures.

\subsubsection{Physiological Scale Boundary}
\label{subsubsec:physiological-scale-boundary}

The present work is posed at the macrocirculatory arterial scale. 
At this scale, blood is commonly modeled as a viscous fluid whose bulk behavior
can be described by continuum mechanics. The goal is not to resolve the effects of individual
blood cells, platelets, or plasma constituents but rather to predict
macroscopic observables such as pressure, flow rate, pressure drop, and
pressure ratio.
\cite[Sec.~1.1.2, pp.~9–10]{KopplHelmig2023DimensionReducedModeling}
\cite[Sec.~1.1.3, pp.~12–14]{KopplHelmig2023DimensionReducedModeling}.

\subsubsection{Blood as a Continuum}
\label{subsubsec:blood-as-a-continuum}

Consider a time-dependent family of spatial fluid domains. For a final time
$T_{\mathrm{final}}>0$, use
$I_T\defined(0,T_{\mathrm{final}})$ for interior-time PDE statements and
$\overline{I}_T\defined[0,T_{\mathrm{final}}]$ for initial data and flow maps.
Define
\begin{flalign*}
    \quad \begin{cases}
     \OmgBt \defined \{\vx \in \R^3 : \vx \text{ lies inside the vessel at time } t\} \\
     \mathcal{Q}_T \defined \{(t,\vx):t\in I_T,\ \vx\in\OmgBt\},
    \end{cases}&&
\end{flalign*}
where the reference configuration is denoted $\OmgB(0)$, the current
configuration is $\OmgBt$, and the space-time fluid domain is $\mathcal{Q}_T$.
This domain notation is the one collected in
Appendix~\ref{app:mathematical-notation}; boundary and function-space
conventions are summarized in Appendix~\ref{app:mathematical-foundations}. In
this report, the continuum assumption is a large-vessel idealization: individual
blood cells, platelets, and plasma constituents are not resolved, and blood
velocity, pressure, density, and Cauchy stress are represented by sufficiently
regular fields on $\mathcal{Q}_T$. Cellular-scale effects may enter only
through later constitutive choices
\cite[Part~I, Ch.~1, Secs.~1.1--1.2, pp.~4--5]{hemodynamics}
\cite[Secs.~1.1.3 and~3.2, pp.~12--14 and~46--47]{KopplHelmig2023DimensionReducedModeling}.
The continuum setup is as follows: Remark~\ref{rmk:eulerian-lagrangian},
Definition~\ref{def:material-derivative},
Definitions~\ref{def:material-density-preservation}--\ref{def:incompressible-flow},
Definitions~\ref{def:linear-momentum-balance}--\ref{def:linear-momentum-balance-strong},
and Definitions~\ref{def:fluidsmodel}--\ref{def:incompressible-stress-tensor}
fix the assumptions used in the model hierarchy.


\subsubsection{Blood Dynamics}
\label{subsubsec:blood-dynamics}

This subsection fixes the kinematic vocabulary used by the later continuum and
reduced-order models.
For $t\in\overline{I}_T$, let
$\La_t:\OmgB(0)\to\OmgB(t)$ be an orientation-preserving
diffeomorphism with $\det\nabla_{\vxi}\La_t>0$.
As a standing assumption, the family $\La_t$ has the time and space regularity
needed to differentiate $\nabla_{\vxi}\La_t$ and
$\det\nabla_{\vxi}\La_t$ in the classical identities below and in
Appendix~\ref{app:continuum-derivation-details}.
It carries each material label $\vxi$ to its current location
$\vx(t,\vxi)=\La_t(\vxi)$.
The Eulerian velocity field $\vu:\mathcal{Q}_T\to\R^3$ satisfies
\begin{flalign*}
    \quad \partial_t\La_t(\vxi)=\vu(t,\La_t(\vxi)). &&
\end{flalign*}
Figure~\ref{fig:lagrangian-eulerian-map} illustrates this Lagrangian--Eulerian
notation.

\begin{figure}[!htb]
    \centering
    \captionsetup{aboveskip=2pt,belowskip=2pt}
    \figtikz{lagrangian-eulerian-map}
    \caption{%
        Lagrangian and Eulerian descriptions. The flow map
        $\La_t:\OmgB(0)\to\OmgBt$ carries the material label
        $\vxi\in\OmgB(0)$ to its current position
        $\vx(t,\vxi)\in\OmgBt$. The Eulerian velocity $\vu(t,\vx)$
        is the velocity of whichever particle currently sits at
        $(t,\vx)$; the Lagrangian velocity is
        $\vuhat(t,\vxi)=\partial_t\vx(t,\vxi)=\vu(t,\La_t(\vxi))$,
        and trajectories $T_{\vxi}$ are generated by the velocity field $\vu$.}
    \label{fig:lagrangian-eulerian-map}
\end{figure}

\begin{remark}[Eulerian vs. Lagrangian]
    \label{rmk:eulerian-lagrangian}
    The Eulerian description views fields such as $\phi(t,\vx)$ or
    $\vu(t,\vx)$ at fixed spatial locations. The Lagrangian description follows
    material particles $\vxi$ along trajectories generated by the velocity field and views fields as functions of the
    material label $\vxi$ and time $t$, such as $\widehat{\phi}(t,\vxi)=\phi(t,\La_t(\vxi))$ or
    $\vuhat(t,\vxi)=\vu(t,\La_t(\vxi))$. The flow map $\La_t$ is the change of
    variables between these descriptions, and the trajectory of a particle with
    label $\vxi$ is the characteristic curve
    $t\mapsto \vx(t;\vxi)=\La_t(\vxi)$. Where defined, the inverse
    $\La_t^{-1}$ maps an Eulerian location back to its material label.
\end{remark}

Under the standard assumption that the prescribed velocity field is continuous in time and uniformly Lipschitz in space, particle trajectories exist uniquely and depend continuously on their initial labels; see \cite[Ch.~2]{Teschl2012OrdinaryDifferentialEquations}. This is a trajectory-level ODE result and should not be confused with Navier--Stokes well-posedness.

\begin{definition}[Material derivative]
    \label{def:material-derivative}
    For $\phi\in C^1(\mathcal{Q}_T)$, the material derivative is the derivative
    of $\phi$ along a particle trajectory:
    \begin{flalign*}
        \quad D_t\phi(t,\vx)
        \defined
        \frac{d}{dt}\phi(t,\La_t(\vxi))
        =
        \partial_t\phi(t,\vx)+\vu(t,\vx)\cdot\nabla\phi(t,\vx),
        \qquad \vxi=\La_t^{-1}(\vx). &&
    \end{flalign*}
    For vector fields, $D_t$ is applied componentwise.
\end{definition}

Let $\vFhat_t=\nabla_{\vxi}\La_t$ be the deformation gradient and
$\Jhatdet_t=\det\vFhat_t$ the Jacobian determinant. In Eulerian notation this
is written $\Jdet_t(\vx)=\det\vFhat_t(\La_t^{-1}(\vx))$. The standard Jacobian evolution law $D_t\Jdet_t=\Jdet_t\,\Div\vu$ is used implicitly below; derivation details are deferred to Appendix~\ref{app:continuum-derivation-details}.

\begin{theorem}[Reynolds Transport Theorem]
    \label{thm:reynolds-transport}
    Let $V(t)=\La_t(V(0))$ be a bounded material volume with Lipschitz boundary.
    If $\phi\in C^1(\mathcal{Q}_T)$, then
    \begin{flalign*}
        \quad \frac{d}{dt}\int_{V(t)}\phi\,\dx
        &=
        \int_{V(t)}\left(D_t\phi+\phi\,\Div\vu\right)\dx \\
        &=
        \int_{V(t)}\left(\partial_t\phi+\Div(\phi\vu)\right)\dx. &&
    \end{flalign*}
\end{theorem}

Proof-level details for Theorem~\ref{thm:reynolds-transport} and the momentum balance below are found
in Appendix~\ref{app:continuum-derivation-details}.

\subsubsection*{Continuum assumptions retained in the main model}

Let the 3D velocity, pressure, and density fields be
\begin{flalign*}
    \quad \vu:\mathcal{Q}_T\to\R^3,\qquad
    p:\mathcal{Q}_T\to\R,\qquad
    \rho:\mathcal{Q}_T\to\R^+. &&
\end{flalign*}
The balance-law vocabulary developed below permits variable density. The reduced arterial models considered later adopt the constant-density specialization $\rho\equiv\rho_0>0$, which is standard for large-vessel incompressible blood-flow models.

The domains and material control volumes are assumed regular enough for traces,
outward unit normals, and the divergence theorem; Lipschitz boundaries are
sufficient for the integral manipulations used here; see
Convention~\ref{conv:foundation-standing-regularity} in
Appendix~\ref{app:mathematical-foundations}.

\begin{definition}[Material density preservation]
    \label{def:material-density-preservation}
    A flow has material density preservation if density is preserved along
    particle trajectories:
    \begin{flalign*}
        \quad D_t\rho=0. &&
    \end{flalign*}
\end{definition}

\begin{definition}[Incompressible flow]
    \label{def:incompressible-flow}
    A flow is incompressible, or volume preserving, if material volumes are
    preserved. With $\La_0=\mathrm{Id}$ this is equivalent to
    \begin{flalign*}
        \quad \Jdet_t\equiv 1
        \quad \Longleftrightarrow \quad
        \Div\vu=0. &&
    \end{flalign*}
\end{definition}

\begin{remark}[Divergence-free condition]
    \label{rmk:incompressibility-divergence-free}
    Mass conservation gives
    \begin{flalign*}
        \quad \partial_t\rho+\Div(\rho\vu)=0
        \quad \Longleftrightarrow \quad
        D_t\rho+\rho\Div\vu=0. &&
    \end{flalign*}
    In the constant-density model used for incompressible Navier--Stokes,
    $\rho\equiv\rho_0>0$, so mass conservation reduces to
    $\Div\vu=0$ in $\mathcal{Q}_T$.
\end{remark}

\begin{definition}[Linear momentum balance]
    \label{def:linear-momentum-balance}
    For a material fluid element $V(t)\subset\OmgBt$, the linear momentum balance is
    \begin{flalign*}
        \quad
        \frac{d}{dt}\int_{V(t)}\rho\vu\,dV
        =
        \int_{\partial V(t)}\vT\nhat\,dS
        +
        \int_{V(t)}\rho\vfb\,dV, &&
    \end{flalign*}
    where $\vT$ is the Cauchy stress tensor and $\vfb$ is body force per unit mass.
\end{definition}

\begin{definition}[Strong momentum balance]
    \label{def:linear-momentum-balance-strong}
    Under the regularity and mass-balance assumptions above, the local momentum
    balance is
    \begin{flalign*}
        \quad
        \rho\left(\partial_t\vu+(\vu\cdot\nabla)\vu\right)
        =
        \Divbf(\vT)+\rho\vfb,
        \qquad \text{in } \OmgBt. &&
    \end{flalign*}
\end{definition}

The stress law closes the momentum equation. Let the rate-of-deformation tensor be
\begin{flalign*}
    \quad \vct{D}(\vu)\defined
    \frac{1}{2}\left(\nabla\vu+\nabla\vu^\top\right). &&
\end{flalign*}

\begin{definition}
    \label{def:newtonian-fluid}
    A fluid is Newtonian if, after separating the pressure contribution, its
    viscous stress depends linearly on the rate-of-deformation tensor
    $\vct{D}(\vu)$.
\end{definition}

\begin{remark}[Blood isotropy justification]
    \label{rmk:blood-isotropy-justification}
    Blood is often modeled as an isotropic continuum fluid at large-vessel
    scales. This is a closure assumption for averaged behavior, not a microscopic
    statement about cells suspended in plasma. The cited sources support the
    scale and rheology caveat--nearly constant large-vessel relative viscosity
    versus diameter- and hematocrit-dependent apparent viscosity in small tubes--not
    a direct cellular-level isotropy claim
    \cite[Sec.~1.1.3, Fig.~1.14, pp.~12--14; Sec.~3.2, pp.~46--47]{KopplHelmig2023DimensionReducedModeling}
    \cite[Abstract; pp.~H1770--H1772]{bviscosity}.
\end{remark}

\begin{remark}[Isotropic constitutive response]
    \label{rmk:isotropy-implication}
    An isotropic fluid has a constitutive response that is independent of the chosen coordinate system. Consequently, the Cauchy stress depends on coordinate-free invariants of $\vct{D}(\vu)$ rather than on a particular basis.
\end{remark}

\begin{definition}[Newtonian, isotropic constitutive law]
    \label{def:fluidsmodel}
    For a Newtonian, isotropic fluid,
    \begin{flalign*}
        \quad \vT=-p\,\ident+2\eta\,\vct{D}(\vu)
        +\lambda\,\Div(\vu)\,\ident, &&
    \end{flalign*}
    where $\eta$ is the dynamic viscosity and $\lambda$ is the coefficient of the
    volumetric-viscous term.
\end{definition}

\begin{definition}[Incompressible stress tensor]
    \label{def:incompressible-stress-tensor}
    When $\Div\vu=0$, the Newtonian incompressible Cauchy stress is
    \begin{flalign*}
        \quad \vT=-p\,\ident+2\eta\,\vct{D}(\vu). &&
    \end{flalign*}
\end{definition}

More general models allow viscosity to vary with shear rate. For generalized Newtonian blood models, the dynamic viscosity is allowed to depend on the local deformation-rate magnitude
\cite[Part~I, Ch.~1, Secs.~2.7--3.1, pp.~16--19]{hemodynamics}.
A standard constitutive form is
\begin{flalign*}
    \quad \dot{\gamma}\defined\sqrt{2\vct{D}(\vu):\vct{D}(\vu)}, &&
\end{flalign*}
\begin{flalign*}
    \quad \vT=-p\,\ident
    +2\eta_{\mathrm{eff}}(\dot{\gamma})\vct{D}(\vu). &&
\end{flalign*}
Here
\begin{flalign*}
    \quad \norm{\vct{D}(\vu)}^2 \defined \vct{D}(\vu):\vct{D}(\vu)
    =\sum_{i,j}D_{ij}(\vu)^2 &&
\end{flalign*}
is the pointwise Frobenius norm squared. When a source writes the viscosity as a
function of \(D:D\) or \(|D|_{\mathrm F}^2\), this report translates it
through the scalar shear-rate convention above.
Here $\eta_{\mathrm{eff}}$ is an effective dynamic viscosity. Kinematic
viscosity is $\nu=\eta/\rho$ after division by density.
The Newtonian incompressible model is recovered from this generalized
Newtonian form when $\eta_{\mathrm{eff}}$ is constant.

\begin{remark}[Newtonian Blood Justification]
    \label{rmk:newtonian-justification}
    For large arteries, empirical studies suggest that apparent viscosity is often
    approximately constant across the relevant operating range, motivating the
    Newtonian approximation in many reduced-order hemodynamic models. However,
    diameter-dependent, hematocrit-dependent, and shear-dependent effects become
    increasingly important in smaller vessels and low-shear regions
    \cite[Sec.~1.1.3, pp.~12–14]{KopplHelmig2023DimensionReducedModeling}
    \cite{bviscosity}.
\end{remark}
